\documentclass[11pt]{article}
\usepackage[utf8]{inputenc}
\usepackage[T1]{fontenc}
\usepackage{graphicx}
\usepackage{amsmath}
\usepackage{amssymb}
\usepackage{hyperref}
\author{Hesam Pakdaman}
\date{\today}
\title{}

\newcommand{\interior}[1]{{\kern0pt#1}^{\mathrm{o}}}

\begin{document}

\begin{itemize}

\item[\noindent\textbf{9.} (a)] If $p \in \interior{E}$ then $p$ is an
  interior point of $E$ which means that there exists some
  neighborhood $N$ with $r > 0$ such that $N \subset E$. By Theorem
  2.19 $N$ is an open set so that every point in $N$ are interior
  points of $E$. This implies that $N \subset \interior{E}$ and
  therefore $\interior{E}$ is an open set.

\item[(b)] If $E$ is open then for every point $x \in E$ we can find a
  neighborhood $N(x)$ with $r > 0$ such that $N \subset E$. This
  implies that every point $x \in E$ is an interior point of $E$ and
  it is therefore true that $x \in \interior{E}$. This shows that
  $E \subset \interior{E}$. Since by construction
  $\interior{E} \subset E$ we have that $\interior{E} = E$.

  If $\interior{E} = E$ then all points in $E$ are interior points of
  $E$. This is clearly true since $\interior{E}$ is the set of all
  interior points of $E$. Then by definition 2.18 (f) $E$ is an open
  set.

\item[(c)] If $G$ is open then we know that for every point $p \in G$
  we can find a neighborhood $N$ with $r > 0$ such that $N \subset
  G$. Since $G \subset E$ we have that $N \subset G \subset E$ which
  shows that $p$ is an interior point of $E$. Hence
  $p \in \interior{E}$ and therefore we have that
  $G \subset \interior{E}$ since $p$ was arbitrary chosen from $G$.

\item[(d)] First we show that
  $(\interior{E})^c \subset \overline{E}^c$. Suppose that
  $x \notin E$, then

  $$ x \in E^c \subset E^c \cup E^{c'} = \overline{E}^c.$$

  Now let $x \in E$. Since $x$ is in the complement to $\interior{E}$
  we know that $x$ is not an interior point to $E$. Therefore for
  every neighborhood $N(x)$ with radius $r > 0$ we have that
  $N \not\subset E$. This means that $N$ always have points in
  $E^c$ which makes $x$ a limit point to $E^c$.  We get that
  $x \in E^{c'} \subset \overline{E}^c$ which shows that
  $(\interior{E})^c \subset \overline{E}^c$. \\

  Conversely, let $p \in \overline{E}^c = E^c \cup E^{c'}$. Then it is
  clear that either $p \in E^c$ or $p \in E^{c'}$. Assume $p \in E^c$,
  then we know that $p$ is not an interior point to $E$ so that
  $p \not\in \interior{E}$ which implies that
  $p \in (\interior{E})^c$. If $p \in E^{c'}$ then it is a limit point
  to $E^c$ and therefore every neighborhood $N(p)$ with radius $r > 0$
  have points (other than $p$) from $E^c$. Hence there is no
  neighborhood such that $N(p) \not\subset E$ which means that $p$
  cannot be an interior point to $E$. We therefore must get that
  $p \in (\interior{E})^c$ and have now shown that
  $\overline{E}^c \subset (\interior{E})^c$. This concludes the proof.

\item[(e)] We shall show that the answer is yes. By (a) we know that
  $\interior{E}$ is always open. Since
  $\interior{E} \subset E \subset \overline{E}$ we can use (c) with
  $G = \interior{E}$ and $E = \overline{E}$ to show that
  $\interior{E} \subset \interior{\overline{E}}$. \\

  Conversely, suppose $x \in \interior{\overline{E}}$ but
  $x \not\in \interior{E}$. This means that $x \in (\interior{E})^c$
  and by (d) we have that

  $$ x \in (\interior{E})^c \Rightarrow x \in \overline{E}^c,$$

  which contradicts our assumption that
  $x \in \interior{\overline{E}}$ since
  $\interior{\overline{E}} \subset \overline{E}$. Hence
  $x \in \interior{\overline{E}}$ implies that $x\in \interior{E}$ and
  we have shown that $\interior{\overline{E}} \subset
  \interior{E}$. This concludes the proof.

\item[(f)] No, and we provide an example. Consider the set
  $E = \{ \frac{1}{n} \}$ with $n = 1, 2, 3, \dots$ in the metric
  space $R^2$. Since $0$ is the only limit point to $E$ we have that
  its closure is $\overline{E} = E \cup \{ 0 \}$. However, none of the
  points in $E$ are interior points and therefore
  $\interior{E} = \emptyset$ which means that
  $\overline{\interior{E}} = \emptyset$. We have shown that
  $\overline{E} \neq \overline{\interior{E}}$ as desired.

\end{itemize}

\hfill $\square$

\end{document}